\documentclass[10pt]{article}

% ---------- Document Identity (update these only) ----------
\newcommand{\DocStage}{}
\newcommand{\DocTitle}{Your Road to Tier One SOF}
\newcommand{\ProgramName}{182-Week Civilian-to-Selection Readiness Pipeline}
\newcommand{\ProgramSubtitle}{}

% ---------- Page ----------
\usepackage[legalpaper,left=0.5in,right=0.5in,top=0.6in,bottom=0.45in]{geometry}

% ---------- Typography ----------
\usepackage[T1]{fontenc}
\usepackage[utf8]{inputenc}
\usepackage{libertinus}
\usepackage{microtype}
\usepackage{parskip}
\usepackage{ragged2e}
\usepackage{textcomp}
\AtBeginDocument{\color{black}}
\AtBeginDocument{\pagecolor{white}}
\AtBeginDocument{\justifying}

% ---------- Landscape pages ----------
\usepackage{pdflscape}

% ---------- Color ----------
\usepackage[table]{xcolor}
\definecolor{StageBlue}{HTML}{1F4E79}
\definecolor{StageBlueDark}{HTML}{163A5A}
\definecolor{LightGray}{HTML}{F2F4F7}
\definecolor{MidGray}{HTML}{D9DEE5}
\definecolor{RuleGray}{HTML}{C7CED8}
\definecolor{HeaderInk}{HTML}{000000}
\definecolor{Ink}{HTML}{000000}

% ---------- Utility ----------
\usepackage{enumitem}
\sloppy
\setlength{\emergencystretch}{2em}

% ---------- Boxes / UI ----------
\usepackage[most]{tcolorbox}

\tcbset{
  charterTitle/.style={
    enhanced,
    colback=StageBlue!4,
    colframe=StageBlue,
    boxrule=1.25pt,
    arc=3pt,
    left=16pt,right=16pt,top=14pt,bottom=14pt,
    borderline north={3.2pt}{0pt}{StageBlue},
    borderline west={4pt}{0pt}{StageBlue},
    borderline south={1.25pt}{0pt}{StageBlue},
    drop shadow={black!25!white},
  },
  charterRule/.style={
    enhanced,
    colback=LightGray,
    colframe=RuleGray,
    boxrule=0.85pt,
    arc=3pt,
    left=12pt,right=12pt,top=10pt,bottom=10pt,
    borderline west={3pt}{0pt}{StageBlue},
  },
  charterBadge/.style={
    enhanced,
    colback=StageBlue,
    coltext=white,
    colframe=StageBlueDark,
    boxrule=0.6pt,
    arc=2pt,
    left=6pt,right=6pt,top=3pt,bottom=3pt,
  }
}
\newtcbox{\Badge}{nobeforeafter,charterBadge}

% ---------- Lists ----------
\setlist[itemize]{leftmargin=1.5em, itemsep=0.25em, topsep=0.25em}
\setlist[enumerate]{leftmargin=1.8em, itemsep=0.25em, topsep=0.25em}

% ---------- TOC ----------
\usepackage{tocloft}
\setcounter{tocdepth}{3}
\setlength{\cftbeforesecskip}{3pt}
\setlength{\cftbeforesubsecskip}{2pt}
\setlength{\cftbeforesubsubsecskip}{0pt}
\renewcommand{\cftsecfont}{\bfseries}
\renewcommand{\cftsecpagefont}{\bfseries}
\renewcommand{\cftsubsecfont}{\normalfont}
\renewcommand{\cftsubsecpagefont}{\normalfont}
\renewcommand{\cftsubsubsecfont}{\normalfont}
\renewcommand{\cftsubsubsecpagefont}{\normalfont}
\setlength{\cftsecindent}{0pt}
\setlength{\cftsubsecindent}{1.25em}
\setlength{\cftsubsubsecindent}{2.8em}
\setlength{\cftsecnumwidth}{2.1em}
\setlength{\cftsubsecnumwidth}{2.8em}
\setlength{\cftsubsubsecnumwidth}{3.6em}

% Remove the built-in TOC heading so only the custom heading is shown
\makeatletter
\renewcommand{\tableofcontents}{\@starttoc{toc}}
\makeatother

% ---------- Header/Footer: page number only ----------
\usepackage{fancyhdr}
\pagestyle{fancy}
\fancyhf{}
\renewcommand{\headrulewidth}{0pt}
\renewcommand{\footrulewidth}{0pt}
\fancyfoot[C]{\thepage}

% ---------- Section formatting ----------
\usepackage{titlesec}

\titleformat{\section}
  {\Large\bfseries\color{Ink}}
  {\thesection}{0.6em}{}
  [\vspace{0.25em}\textcolor{StageBlue}{\rule{\linewidth}{0.8pt}}]

\titleformat{\subsection}
  {\large\bfseries\color{Ink}}
  {\thesubsection}{0.6em}{}

\titleformat{\subsubsection}
  {\normalsize\bfseries\color{Ink}}
  {\thesubsubsection}{0.6em}{}

\titlespacing*{\section}{0pt}{0.95em}{0.35em}
\titlespacing*{\subsection}{0pt}{0.65em}{0.25em}
\titlespacing*{\subsubsection}{0pt}{0.55em}{0.2em}

% ---------- Table engine ----------
\usepackage{tabularray}
\UseTblrLibrary{booktabs}

% ---------- tabularray: GLOBAL table treatment ----------
\SetTblrInner{
  cells = {fg=black, bg=white, valign=t},
}

% ---------- Header helpers ----------
\newcommand{\Hdr}[1]{\shortstack[c]{\strut #1\strut}}

% ---------- Lines ----------
\newcommand{\TitleLine}{\textcolor{StageBlue}{\rule{0.98\linewidth}{0.95pt}}}
\newcommand{\SubTitleLine}{\textcolor{RuleGray}{\rule{0.98\linewidth}{0.65pt}}}

% ---------- Continued on next page ----------
\newcommand{\ContinuedOnNextPageInline}{%
  \par
  \vfill
  \noindent\textcolor{RuleGray}{\rule{\linewidth}{1pt}}\vspace{-2pt}
  \noindent\hfill\textit{Continued on next page}\par\vspace{-10pt}
  \noindent\textcolor{RuleGray}{\rule{\linewidth}{1pt}}
  \clearpage
}

% ---------- PDF hyperlinks ----------
\usepackage[hidelinks]{hyperref}
\hypersetup{
  colorlinks=false,
  pdfborder={0 0 0},
  linktoc=all,
  pdftitle={\DocTitle},
  pdfauthor={\ProgramName}
}

% =========================================================
\begin{document}



\section*{Stage 1.E — Standardized Templates}
\phantomsection
\addcontentsline{toc}{section}{Stage 1.E — Standardized Templates}

\subsection*{1.E.1 Overview Template}
\phantomsection
\addcontentsline{toc}{subsection}{1.E.1 Overview Template}

\begin{verbatim}
Overview \textbf{ID}: [\textbf{REQUIRED}: 7.01–7.20]
Overview Name: [\textbf{REQUIRED}]
Version: [\textbf{REQUIRED}: vX.Y (baseline v1.0 unless patched)]
Governing Status: [\textbf{REQUIRED}: Draft — Non-Deployable | Deployable]
Scope Boundary Statement: [\textbf{REQUIRED}: What this covers / What this does not cover]

Owner (primary executor): [\textbf{OPTIONAL}: Coach Lead | Trainee | Medical | Equipment Lead]
Patch References: [\textbf{OPTIONAL}: \textbf{PATCH}-1E-### if schema changes apply, plus any
artifact-specific patch IDs if used elsewhere]
\end{verbatim}

\subsubsection*{1.E.1.2 Overview Required Sections}
\phantomsection

Rule: The following headings are template-locked. Do not rename. Populate all required fields to meet minimum completion standards. Do not embed programming prescriptions, test protocols, calendars, or shopping lists.

1) Purpose \& Scope

\begin{itemize}
\item Purpose: [\textbf{REQUIRED}: One paragraph describing the capability domain covered and why it exists in this program.]
\item Scope boundaries: [\textbf{REQUIRED}: In-scope and out-of-scope bullets; must include procurement scope and sourcing boundaries if procurement is discussed.]
\item Crosswalk intent: [\textbf{REQUIRED}: One sentence stating how this overview ties into phases and evidence posture.]
\item Prohibited-content boundary: [\textbf{REQUIRED}: One bullet explicitly stating what this overview does NOT provide (no workouts, no test protocols, no calendars, no medical advice).]
\end{itemize}

2) Governance And Safety Boundaries

\begin{itemize}
\item Boundary alignment: [\textbf{REQUIRED}: State relevant Stage 0 boundaries applicable to this overview.]
\item Prohibitions (if applicable): [\textbf{REQUIRED}: Any \textbf{MUST} NOT items relevant to this domain, including OTC-only posture, no prescription guidance, no hormone guidance, no illegal procurement guidance.]
\item Water governance (if applicable): [\textbf{REQUIRED}: Surface-only default unless explicitly authorized later; no unsupervised underwater/breath-hold ever; certified supervision required for any underwater/breath-hold exposure.]
\item Escalation posture: [\textbf{REQUIRED}: Red-flag triggers and “stop/escalate when indicated” posture; no medical advice.]
\item Evidence posture: [\textbf{REQUIRED}: Statement that unsafe or boundary-violating behavior produces inadmissible evidence and cannot support readiness claims.]
\end{itemize}

3) Tier Doctrine

\begin{itemize}
\item Tier definitions: [\textbf{REQUIRED}: Tier 0–Tier 3 described (see Field Definitions).]
\item Price band convention: [\textbf{REQUIRED}: Use \$ / \$\$ / \$\$\$ only; definition included below.]
\item N/A rule: [\textbf{REQUIRED}: Declare when N/A may be used (only where explicitly permitted; see Field Definitions).]
\item Substitution posture: [\textbf{REQUIRED}: Substitutions preserve intent and reduce risk; explicitly state what breaks intent.]
\item Minimal/Standard/Optimal mapping (optional): [\textbf{OPTIONAL}: If included, must map to Tier 0–Tier 3 without renumbering tiers.]
\end{itemize}

4) Tiered Kit (Tables)

Rule: Tables are required for Tier 0–Tier 3. Use category/spec language or explicit brand references as needed. No illegal sourcing workarounds.

Tier 0 — Baseline (Start-from-Zero / Household-Only / None)

\begingroup
\scriptsize
\setlength{\tabcolsep}{2pt}
\renewcommand{\arraystretch}{1.12}
\begin{longtblr}{
  width=\linewidth,
  rowhead=1,
  colspec = {
    X[l,1.10]
    X[l,1.25]
    X[l,1.35]
    Q[l,wd=1.05in]
    Q[l,wd=0.95in]
    X[l,1.25]
    Q[l,wd=0.95in]
    X[l,1.55]
  },
  hlines,
  vlines,
  cells = {hsep=6pt, vsep=6pt, halign=l, valign=t},
  row{odd} = {bg=white},
  row{even} = {bg=LightGray},
  row{1} = {bg=StageBlue!15, fg=black, font=\bfseries, halign=l, valign=t},
}
\Hdr{Item} & \Hdr{Purpose} & \Hdr{Key Specs} & \Hdr{Typical Footprint} & \Hdr{Approx. Price Band} & \Hdr{Maintenance} & \Hdr{Phase Tie-In} & \Hdr{Notes} \\
{[\textbf{REQUIRED}]} & {[\textbf{REQUIRED}]} & {[\textbf{REQUIRED}]} & {[\textbf{REQUIRED}]} & {[\$ / \$\$ / \$\$\$]} & {[\textbf{REQUIRED}]} & {[\textbf{REQUIRED}]} & {[\textbf{REQUIRED}]} \\
\end{longtblr}
\endgroup

Tier 1 Essentials

\begingroup
\scriptsize
\setlength{\tabcolsep}{2pt}
\renewcommand{\arraystretch}{1.12}
\begin{longtblr}{
  width=\linewidth,
  rowhead=1,
  colspec = {
    X[l,1.10]
    X[l,1.25]
    X[l,1.35]
    Q[l,wd=1.05in]
    Q[l,wd=0.95in]
    X[l,1.25]
    Q[l,wd=0.95in]
    X[l,1.55]
  },
  hlines,
  vlines,
  cells = {hsep=6pt, vsep=6pt, halign=l, valign=t},
  row{odd} = {bg=white},
  row{even} = {bg=LightGray},
  row{1} = {bg=StageBlue!15, fg=black, font=\bfseries, halign=l, valign=t},
}
\Hdr{Item} & \Hdr{Purpose} & \Hdr{Key Specs} & \Hdr{Typical Footprint} & \Hdr{Approx. Price Band} & \Hdr{Maintenance} & \Hdr{Phase Tie-In} & \Hdr{Notes} \\
{[\textbf{REQUIRED}]} & {[\textbf{REQUIRED}]} & {[\textbf{REQUIRED}]} & {[\textbf{REQUIRED}]} & {[\$ / \$\$ / \$\$\$]} & {[\textbf{REQUIRED}]} & {[\textbf{REQUIRED}]} & {[\textbf{REQUIRED}]} \\
\end{longtblr}
\endgroup

Tier 2 Upgrades

\begingroup
\scriptsize
\setlength{\tabcolsep}{2pt}
\renewcommand{\arraystretch}{1.12}
\begin{longtblr}{
  width=\linewidth,
  rowhead=1,
  colspec = {
    X[l,1.10]
    X[l,1.25]
    X[l,1.35]
    Q[l,wd=1.05in]
    Q[l,wd=0.95in]
    X[l,1.25]
    Q[l,wd=0.95in]
    X[l,1.55]
  },
  hlines,
  vlines,
  cells = {hsep=6pt, vsep=6pt, halign=l, valign=t},
  row{odd} = {bg=white},
  row{even} = {bg=LightGray},
  row{1} = {bg=StageBlue!15, fg=black, font=\bfseries, halign=l, valign=t},
}
\Hdr{Item} & \Hdr{Purpose} & \Hdr{Key Specs} & \Hdr{Typical Footprint} & \Hdr{Approx. Price Band} & \Hdr{Maintenance} & \Hdr{Phase Tie-In} & \Hdr{Notes} \\
{[\textbf{REQUIRED}]} & {[\textbf{REQUIRED}]} & {[\textbf{REQUIRED}]} & {[\textbf{REQUIRED}]} & {[\$ / \$\$ / \$\$\$]} & {[\textbf{REQUIRED}]} & {[\textbf{REQUIRED}]} & {[\textbf{REQUIRED}]} \\
\end{longtblr}
\endgroup

Tier 3 Advanced

\begingroup
\scriptsize
\setlength{\tabcolsep}{2pt}
\renewcommand{\arraystretch}{1.12}
\begin{longtblr}{
  width=\linewidth,
  rowhead=1,
  colspec = {
    X[l,1.10]
    X[l,1.25]
    X[l,1.35]
    Q[l,wd=1.05in]
    Q[l,wd=0.95in]
    X[l,1.25]
    Q[l,wd=0.95in]
    X[l,1.55]
  },
  hlines,
  vlines,
  cells = {hsep=6pt, vsep=6pt, halign=l, valign=t},
  row{odd} = {bg=white},
  row{even} = {bg=LightGray},
  row{1} = {bg=StageBlue!15, fg=black, font=\bfseries, halign=l, valign=t},
}
\Hdr{Item} & \Hdr{Purpose} & \Hdr{Key Specs} & \Hdr{Typical Footprint} & \Hdr{Approx. Price Band} & \Hdr{Maintenance} & \Hdr{Phase Tie-In} & \Hdr{Notes} \\
{[\textbf{REQUIRED}]} & {[\textbf{REQUIRED}]} & {[\textbf{REQUIRED}]} & {[\textbf{REQUIRED}]} & {[\$ / \$\$ / \$\$\$]} & {[\textbf{REQUIRED}]} & {[\textbf{REQUIRED}]} & {[\textbf{REQUIRED}]} \\
\end{longtblr}
\endgroup

5) Selection Criteria \& Fit Checks

\begin{itemize}
\item Selection criteria: [\textbf{REQUIRED}: What functional requirements matter (capacity, fit, durability, safety, measurability, compatibility). No illegal sourcing guidance.]
\item Fit checks: [\textbf{REQUIRED}: Pass/fail checks a coach or trainee can perform (sizing, stability, comfort, compatibility, measurability).]
\item Common mismatch risks: [\textbf{REQUIRED}: How wrong selection causes safety/validity issues (e.g., foot breakdown, inaccurate timing, unsafe anchoring, inconsistent measurements).]
\item Documentation tie-in: [\textbf{REQUIRED}: What must be recorded to prove fit checks were completed and to support auditability.]
\end{itemize}

6) Setup, Safety, And Anchoring

\begin{itemize}
\item Setup requirements: [\textbf{REQUIRED}: How to set up safely (space, anchoring, stability, inspection).]
\item Safety controls: [\textbf{REQUIRED}: What \textbf{MUST} be checked before use; include “do not use if” conditions.]
\item Abort posture: [\textbf{REQUIRED}: Stop/abort triggers relevant to this equipment/system.]
\item Environment constraints: [\textbf{REQUIRED}: Any conditions that require substitution (ice, low visibility, poor air quality, facility rules).]
\end{itemize}

7) Maintenance, Replacement, And Storage

\begin{itemize}
\item Maintenance routine: [\textbf{REQUIRED}: Cleaning, inspection, drying, and storage posture.]
\item Replacement triggers: [\textbf{REQUIRED}: Observable wear/failure criteria; specify “replace when” cues.]
\item Storage constraints: [\textbf{REQUIRED}: Apartment/limited space considerations; include drying/hygiene constraints when relevant.]
\item Sustainment risk: [\textbf{REQUIRED}: What predictable failures occur if maintenance/storage is not executed.]
\end{itemize}

8) Substitution Rules (What Is Acceptable; What Breaks Intent)

\begin{itemize}
\item Acceptable substitutions: [\textbf{REQUIRED}: Lower-risk equivalents; what preserves intent.]
\item Unacceptable substitutions: [\textbf{REQUIRED}: What breaks intent, increases risk, violates boundaries, or invalidates evidence.]
\item Documentation requirement: [\textbf{REQUIRED}: What must be logged when substitutions occur (what changed, why, risk posture, whether evidence admissibility is preserved or degraded).]
\item Phase gating: [\textbf{REQUIRED}: If substitutions vary by phase, state the phase gate logic without prescribing workouts.]
\end{itemize}

9) Phase Integration Map (Phase 00–13)

\begin{itemize}
\item Mapping rule: [\textbf{REQUIRED}: Each phase must be referenced at least once; use the table below or equivalent consistent format.]
\item Status definitions: [\textbf{REQUIRED}: Use Introduced / Optional / Required / Not Used only; define them in Tier Doctrine or Field Definitions.]
\end{itemize}

\begingroup
\scriptsize
\setlength{\tabcolsep}{2pt}
\renewcommand{\arraystretch}{1.12}
\begin{longtblr}{
  width=\linewidth,
  rowhead=1,
  colspec = {
    Q[l,wd=1.0in]
    Q[l,wd=2.0in]
    X[l,3.10]
  },
  hlines,
  vlines,
  cells = {hsep=6pt, vsep=6pt, halign=l, valign=t},
  row{odd} = {bg=white},
  row{even} = {bg=LightGray},
  row{1} = {bg=StageBlue!15, fg=black, font=\bfseries, halign=l, valign=t},
}
\Hdr{Phase} & \Hdr{Status} & \Hdr{Notes} \\
Phase 00 & {[Introduced / Optional / Required / Not Used]} & {[\textbf{REQUIRED}]} \\
Phase 01 & {[Introduced / Optional / Required / Not Used]} & {[\textbf{REQUIRED}]} \\
Phase 02 & {[Introduced / Optional / Required / Not Used]} & {[\textbf{REQUIRED}]} \\
Phase 03 & {[Introduced / Optional / Required / Not Used]} & {[\textbf{REQUIRED}]} \\
Phase 04 & {[Introduced / Optional / Required / Not Used]} & {[\textbf{REQUIRED}]} \\
Phase 05 & {[Introduced / Optional / Required / Not Used]} & {[\textbf{REQUIRED}]} \\
Phase 06 & {[Introduced / Optional / Required / Not Used]} & {[\textbf{REQUIRED}]} \\
Phase 07 & {[Introduced / Optional / Required / Not Used]} & {[\textbf{REQUIRED}]} \\
Phase 08 & {[Introduced / Optional / Required / Not Used]} & {[\textbf{REQUIRED}]} \\
Phase 09 & {[Introduced / Optional / Required / Not Used]} & {[\textbf{REQUIRED}]} \\
Phase 10 & {[Introduced / Optional / Required / Not Used]} & {[\textbf{REQUIRED}]} \\
Phase 11 & {[Introduced / Optional / Required / Not Used]} & {[\textbf{REQUIRED}]} \\
Phase 12 & {[Introduced / Optional / Required / Not Used]} & {[\textbf{REQUIRED}]} \\
Phase 13 & {[Introduced / Optional / Required / Not Used]} & {[\textbf{REQUIRED}]} \\
\end{longtblr}
\endgroup

10) Risks, Contraindications, And Red Flags

\begin{itemize}
\item Risks: [\textbf{REQUIRED}: List practical risks in program terms; may reference \textbf{RISK-1B}-\#\#\#.]
\item Contraindications: [\textbf{REQUIRED}: Operational contraindications (do not diagnose); specify “prohibit exposure” posture when relevant.]
\item Red flags: [\textbf{REQUIRED}: Observable triggers requiring stop/escalation.]
\item Gate posture: [\textbf{REQUIRED}: Default posture when risk flags appear (Hold/Regress/Medical Hold as applicable).]
\item Evidence posture: [\textbf{REQUIRED}: Statement that evidence produced under red-flag continuation is inadmissible.]
\end{itemize}

11) Compliance Checklist (Pass/Fail Bullets)

\begin{itemize}
\item {[\textbf{REQUIRED}: Pass/fail checklist; must match 1.E.1.4 requirements at minimum.]}
\item {[\textbf{REQUIRED}: Must include an explicit fail item for any missing required heading or missing Tier table.]}
\end{itemize}

\subsubsection*{1.E.1.3 Field Definitions}
\phantomsection

A) Header Fields (Overview Header Block)

Overview \textbf{ID} (Required)

\begin{itemize}
\item Allowed values: 7.01 through 7.20 only.
\item Completion rule: must match an allowed value exactly (including leading zero). Must remain stable for crosswalks.
\item Status impact if missing/invalid: Governing Status \textbf{MUST} be Draft — Non-Deployable.
\end{itemize}

Overview Name (Required)

\begin{itemize}
\item Allowed values: free text.
\item Completion rule: must describe the capability domain unambiguously; must remain stable for crosswalk use.
\item Status impact if missing: Draft — Non-Deployable.
\end{itemize}

Version (Required)

\begin{itemize}
\item Allowed values: vX.Y where X and Y are integers. Baseline is v1.0 unless patched.
\item Completion rule: version must match the binder’s version state; any change in meaning requires a patch posture and version increment.
\item Status impact if missing/invalid: Draft — Non-Deployable.
\end{itemize}

Governing Status (Required; locked strings)

\begin{itemize}
\item Allowed values only: Draft — Non-Deployable \,|\, Deployable.
\item Definition (must be used verbatim in meaning): Deployable means the artifact’s required fields are complete and its applicable gate checklist has been passed and recorded. Do not append additional label text.
\item Completion rule: set to Draft — Non-Deployable by default; set to Deployable only after the overview passes its compliance checklist and any applicable gate pass record is present.
\item Status impact if missing/invalid: Draft — Non-Deployable.
\end{itemize}

Scope Boundary Statement (Required)

\begin{itemize}
\item Must include, at minimum:
  \begin{itemize}
  \item what the overview covers,
  \item what it does not cover,
  \item explicit prohibition posture for prohibited content categories relevant to the domain (e.g., no medical advice; no prescription or hormone guidance; no illegal procurement guidance; procurement scope and sourcing boundaries; no test protocols).
  \end{itemize}
\item Completion rule: must be written so a reviewer can determine scope without reading the full document.
\item Status impact if missing: Draft — Non-Deployable.
\end{itemize}

Owner (primary executor) (Optional)

\begin{itemize}
\item Allowed values: Coach Lead \,|\, Trainee \,|\, Medical \,|\, Equipment Lead.
\item Completion rule: if populated, must match exactly one allowed value.
\item Status impact if missing: none.
\end{itemize}

Patch References (Optional)

\begin{itemize}
\item Allowed values: \textbf{PATCH}-1E-\#\#\# only for Stage 1.E schema changes; other patch references allowed only if they exist elsewhere and are clearly scoped.
\item Completion rule: include only when (a) the template schema changed and (b) the instantiated overview conforms to the patched schema.
\item Status impact if incorrect or misleading: Draft — Non-Deployable until corrected.
\end{itemize}

B) Governing Status Semantics (Applicable to All Instantiated Overviews)

Draft — Non-Deployable means at least one of the following is true:

\begin{itemize}
\item any required header field is missing/invalid,
\item any required section heading is missing/renamed,
\item any required table is missing or column headers are altered,
\item any required field is left empty beyond allowed placeholders,
\item prohibited content appears,
\item the compliance checklist is not passed.
\end{itemize}

Deployable means all of the following are true:

\begin{itemize}
\item all required header fields present and valid,
\item all required sections present and complete to minimum standards,
\item all required tier tables present with exact column headers,
\item the compliance checklist passes and the pass record exists.
\end{itemize}

C) Tier Doctrine (Tier 0–Tier 3) (Binding Rules)

Tier numbering is canonical. Tier labels must remain: Tier 0, Tier 1, Tier 2, Tier 3.

Tier 0 (Required table section)

\begin{itemize}
\item Meaning: start-from-zero / household-only / none-assumed baseline.
\item Completion rule: Tier 0 \textbf{MUST} be represented as a table and \textbf{MUST} include at least one row.
\item Allowed Tier 0 “none” handling: Tier 0 may include a row where Item explicitly states “None / household-only baseline,” but the row must still include Purpose, Key Specs, Typical Footprint, Maintenance/Consumables, Phase Tie-In, and Notes.
\item Approx. Price Band in Tier 0: N/A may be used only when the tier truly contains no acquisition items; otherwise use \$ / \$\$ / \$\$\$.
\end{itemize}

Tier 1 (Essentials)

\begin{itemize}
\item Meaning: minimum capability categories required to execute safely and maintain validity.
\item Completion rule: entries must be category/spec based; no brand names, no retailer references, no links.
\end{itemize}

Tier 2 (Upgrades)

\begin{itemize}
\item Meaning: improves sustainability, safety margin, or efficiency without changing boundaries.
\item Completion rule: may indicate mandatory procurement when governance or phase execution requires it; must remain compatible with start-from-zero posture.
\end{itemize}

Tier 3 (Advanced)

\begin{itemize}
\item Meaning: advanced capabilities; may increase robustness but must not introduce prohibited practices.
\item Completion rule: must not be treated as required for minimal execution unless explicitly marked Required in Phase Integration Map with justification.
\end{itemize}

Minimal / Standard / Optimal mapping (Optional)

\begin{itemize}
\item If included, it must be an explanatory mapping only and must not replace Tier 0–Tier 3.
\end{itemize}

D) Price Band Convention (Binding)

Approx. Price Band (Required column)

\begin{itemize}
\item Allowed values: \$ \,|\, \$\$ \,|\, \$\$\$ only (plus N/A only in Tier 0 where truly none).
\item Semantics: ordinal category only, internal to the program; not a quote.
\item Numeric ranges: optional. If used, they must be declared once inside the overview (Tier Doctrine) and treated as frozen unless changed via \textbf{PATCH}-1E-\#\#\#.
\end{itemize}

E) Tiered Kit Table Column Definitions (Exact Columns; Completion Rules)

Item

\begin{itemize}
\item Meaning: generic capability descriptor (category-level).
\item Allowed: descriptive text, including brand references when intentionally used.
\item Complete when: a reader can identify the capability class without a specific product.
\end{itemize}

Purpose

\begin{itemize}
\item Meaning: why the item exists (safety, validity, feasibility, maintenance, measurement, capability enablement).
\item Complete when: purpose ties to program constraints, validity, safety, or feasibility.
\end{itemize}

Key Specs

\begin{itemize}
\item Meaning: the minimum functional specification(s) needed to meet purpose.
\item Allowed: measurable specs and interface requirements; prohibited: brand-specific model numbers.
\item Complete when: a coach/trainee can verify compliance via fit checks.
\end{itemize}

Typical Footprint

\begin{itemize}
\item Meaning: storage/space burden for constrained living environment.
\item Allowed: at least one of the following must appear: Pocket | Handheld | Small | Medium | Large | Floor-space | Fixed.
\item Optional: add approximate dimensions in U.S. customary and/or metric.
\item Complete when: storage feasibility is obvious.
\end{itemize}

Approx. Price Band

\begin{itemize}
\item Allowed: \$ / \$\$ / \$\$\$ only (Tier 0: N/A only if truly none).
\item Complete when: exactly one allowed symbol is used.
\end{itemize}

Maintenance/Consumables

\begin{itemize}
\item Meaning: sustainment needs (cleaning, inspection cadence, consumables categories, replacement intervals).
\item Allowed: category-level; prohibited: medical claims.
\item Complete when: a reader can anticipate ongoing burden.
\end{itemize}

Phase Tie-In

\begin{itemize}
\item Meaning: tie the item to phase relevance without calendars.
\item Required format: must reference phases using Phase 00 through Phase 13.
\item Allowed statuses for Phase Tie-In: Introduced | Optional | Required | Not Used.
\item Complete when: the earliest “Introduced” and any “Required” phases are explicit.
\end{itemize}

Notes

\begin{itemize}
\item Meaning: constraints, gating, compatibility, substitution notes, known risks, audit notes.
\item Complete when: any special constraints that would affect safe use or admissibility are captured.
\end{itemize}

F) Phase Integration Map Rules (Binding)

\begin{itemize}
\item Must include Phase 00 through Phase 13 explicitly.
\item Allowed Status values: Introduced | Optional | Required | Not Used only.
\item Notes field must explain the status choice when non-obvious (e.g., why a tier is Required later).
\item Must not include week numbers, dates, deload placement, or test calendar logic.
\end{itemize}

G) Required Section Completion Standards (Minimum)

For each of the 11 required headings:

\begin{itemize}
\item Headings must exist exactly as written.
\item Each heading must contain content beyond placeholders sufficient to meet minimum definitions below.
\item Empty headings fail compliance and force Draft — Non-Deployable.
\end{itemize}

1) Purpose \& Scope

\begin{itemize}
\item Minimum: purpose paragraph; in-scope/out-of-scope bullets; crosswalk intent; prohibited-content boundary.
\item Prohibited: programming, protocols, calendars, medical guidance, procurement links.
\end{itemize}

2) Governance And Safety Boundaries

\begin{itemize}
\item Minimum: applicable boundaries; domain-specific \textbf{MUST} NOT items; escalation posture; evidence inadmissibility statement.
\item Prohibited: diagnosis/treatment, prescription/hormone instructions, underwater/breath-hold permissions without certified supervision.
\end{itemize}

3) Tier Doctrine

\begin{itemize}
\item Minimum: Tier 0–Tier 3 definitions; price band convention; N/A rule; substitution posture.
\item Prohibited: renumbering tiers; converting tiers into shopping lists.
\end{itemize}

4) Tiered Kit (Tables)

\begin{itemize}
\item Minimum: Tier 0–Tier 3 tables present; at least one row in each; exact column headers used.
\item Prohibited: illegal sourcing workarounds and unphased required-purchase language without governance justification.
\end{itemize}

5) Selection Criteria \& Fit Checks

\begin{itemize}
\item Minimum: selection criteria; pass/fail fit checks; mismatch risks; documentation tie-in.
\item Prohibited: subjective-only criteria without a pass/fail check.
\end{itemize}

6) Setup, Safety, And Anchoring

\begin{itemize}
\item Minimum: setup requirements; pre-use checks; abort posture; environment constraints.
\item Prohibited: unsafe improvisation or bypassing facility rules.
\end{itemize}

7) Maintenance, Replacement, And Storage

\begin{itemize}
\item Minimum: maintenance routine; replacement triggers; storage constraints; sustainment risk.
\item Prohibited: unsafe chemical handling guidance; medical claims.
\end{itemize}

8) Substitution Rules (What Is Acceptable; What Breaks Intent)

\begin{itemize}
\item Minimum: acceptable substitutions; unacceptable substitutions; documentation requirement; phase gating notes if applicable.
\item Prohibited: risk-increasing substitutions; make-up volume or intensity behavior; undocumented substitutions.
\end{itemize}

9) Phase Integration Map (Phase 00–13)

\begin{itemize}
\item Minimum: all phases listed; statuses selected from allowed set; notes included.
\item Prohibited: calendars or dates.
\end{itemize}

10) Risks, Contraindications, And Red Flags

\begin{itemize}
\item Minimum: risks; contraindications (operational); red flags; gate posture; evidence posture.
\item Prohibited: medical diagnosis/treatment, medication adjustment advice.
\end{itemize}

11) Compliance Checklist (Pass/Fail Bullets)

\begin{itemize}
\item Minimum: explicit pass conditions and explicit fail conditions; must include missing-heading and missing-table fails.
\item Prohibited: non-auditable “best effort” language.
\end{itemize}

H) Conservative Defaults and Placeholder Rules (Binding)

\begin{itemize}
\item Placeholders are permitted only while an overview remains Draft — Non-Deployable.
\item Deployable requires all required fields to be populated with non-placeholder content meeting completion criteria.
\item Allowed placeholder tokens (Draft only): [\textbf{REQUIRED}], [\textbf{OPTIONAL}], [\textbf{TBD}], N/A (only where explicitly allowed).
\item If any required field remains placeholder at the time of attempted deployable labeling, status \textbf{MUST} remain Draft — Non-Deployable.
\end{itemize}

\subsubsection*{1.E.1.4 Overview Template Compliance Checklist}
\phantomsection

\textbf{PASS} only if every item below is true:

\begin{itemize}
\item Overview \textbf{ID} is correctly formatted 7.01–7.20 (leading zero required) and Overview Name is present.
\item Version is present and formatted vX.Y.
\item Governing Status is set to an allowed value (Draft — Non-Deployable or Deployable) and is consistent with checklist outcome.
\item Scope Boundary Statement is present and includes both in-scope and out-of-scope constraints, including procurement scope and sourcing boundaries when procurement is discussed.
\item All 11 required headings are present exactly as written (no renaming) and contain non-empty content meeting minimum standards.
\item Tier 0, Tier 1, Tier 2, Tier 3 tables exist and use exact column headers:
  Item | Purpose | Key Specs | Typical Footprint | Approx. Price Band | Maintenance/Consumables | Phase Tie-In | Notes
\item Tier 0 is a table section (not narrative only) and contains at least one row.
\item Approx. Price Band uses only \$ / \$\$ / \$\$\$ in Tier 1–Tier 3, and uses only \$ / \$\$ / \$\$\$ or N/A (Tier 0 only) in Tier 0.
\item Phase Integration Map exists and references Phase 00 through Phase 13 using only allowed statuses (Introduced / Optional / Required / Not Used).
\item Substitution Rules section explicitly states acceptable substitutions, unacceptable substitutions, and documentation requirements.
\item Governance And Safety Boundaries section includes all applicable Stage 0 boundary alignment; if water-relevant, includes certified supervision requirement and prohibition on unsupervised underwater/breath-hold exposure.
\item Risks, Contraindications, And Red Flags section includes observable red flags and default gate posture statements (Hold/Regress/Medical Hold as applicable).
\item No prohibited content appears (no workouts, no test protocols, no calendars, no prescription/hormone guidance, no illegal procurement guidance).
\end{itemize}

\textbf{FAIL} if any item above is not satisfied:

\begin{itemize}
\item If any required item fails → Governing Status \textbf{MUST} be Draft — Non-Deployable.
\end{itemize}

\newpage
\subsection*{1.E.2 Phase Row Template}
\phantomsection
\addcontentsline{toc}{subsection}{1.E.2 Phase Row Template}

\subsubsection*{1.E.2.1 Phase Header Block}
\phantomsection

\begin{verbatim}
Phase \textbf{ID}: [\textbf{REQUIRED}: 00–13]
Phase Name: [\textbf{REQUIRED}]
Duration (weeks): [\textbf{REQUIRED}]
Version: [\textbf{REQUIRED}: vX.Y (baseline v1.0 unless patched)]
Governing Status: [\textbf{REQUIRED}: Draft — Non-Deployable | Deployable]

Owner (primary executor): [\textbf{OPTIONAL}: Coach Lead | Trainee | Medical | Equipment Lead]
Patch References: [\textbf{OPTIONAL}: \textbf{PATCH}-1E-### if schema changes apply, plus any artifact-specific patch IDs if used elsewhere]
\end{verbatim}

\subsubsection*{1.E.2.2 Phase Required Sections}
\phantomsection

Rule: The following headings are template-locked. Do not rename. Populate all required fields to meet minimum completion standards. Do not embed training prescriptions (sets/reps/intervals), test protocol cards, or calendars.

1) Phase Mission

\begin{itemize}
\item {[\textbf{REQUIRED}: One paragraph defining what the phase must accomplish in operational terms.]}
\item {[\textbf{REQUIRED}: One sentence stating what “success looks like” for this phase in admissible evidence terms.]}
\end{itemize}

2) Phase Purpose (Why This Phase Exists)

\begin{itemize}
\item {[\textbf{REQUIRED}: Why this phase exists in the continuum; what problem it solves; what it does not attempt.]}
\item {[\textbf{REQUIRED}: Phase boundary statement: what is out-of-scope for this phase (prevents drift).]}
\end{itemize}

3) Entry Criteria (Gates To Start)

\begin{itemize}
\item {[\textbf{REQUIRED}: Objective entry criteria; must reference dependencies by \textbf{ID} when applicable (\textbf{DEP-1C}-\#\#\# / \textbf{ASM-1C}-\#\#\#).]}
\item {[\textbf{REQUIRED}: Required evidence posture: what must be demonstrated using admissible logs/tests (IDs only).]}
\item {[\textbf{REQUIRED}: Fail posture: what happens if entry criteria are not met (Substitute / Delay escalation / \textbf{HOLD} / \textbf{REGRESS} / \textbf{MEDICAL} \textbf{HOLD} as applicable).]}
\item {[\textbf{REQUIRED}: If medical clearance or restrictions are relevant, state “clinician-led clearance when indicated” posture without medical advice.]}
\end{itemize}

4) Operating Constraints

This section must restate, verbatim in meaning (not necessarily word-for-word), the operating constraints that govern execution:

\begin{itemize}
\item Cadence: 6 sessions/week.
\item Session elements: Safety self-check, Warm-up, Mobility, Main work, Cardio, Cooldown, Recovery, Post-session notes.
\item Cardio minimum standard (in body text): Cardio must include a minimum 30 minutes continuous per training session; if Cardio is satisfied within main work and meets the minimum, no supplemental Cardio is required.
\item Steps baseline: 10,000 steps/day baseline (daily; not a session).
\item Safety override: safety overrides schedule and performance.
\item Validity posture: evidence must be admissible; gate decisions rely on \textbf{VALID} evidence only (definitions exist elsewhere; reference posture only here).
\item No stacking: no two-a-days; no make-up doubling behavior.
\item Documentation discipline: same-day logging required; missing required fields make evidence inadmissible for gate use.
\end{itemize}

5) Primary Adaptations And Physiological Focus

\begin{itemize}
\item {[\textbf{REQUIRED}: Bullet list of adaptation targets; no programming detail.]}
\item {[\textbf{REQUIRED}: Identify at least one “protected constraint” for the phase (e.g., tissue tolerance first, impact-managed, load-managed) without prescriptions.]}
\end{itemize}

6) Primary Modalities (Strength/Conditioning/Ruck/Swim/Etc.)

\begin{itemize}
\item {[\textbf{REQUIRED}: List allowed modality domains; note higher-risk exposures that require governed introduction (impact running, loaded rucking, aquatic/underwater governance).]}
\item {[\textbf{REQUIRED}: Prohibited modalities or prohibited conditions for this phase (e.g., underwater/breath-hold unsupervised always prohibited).]}
\end{itemize}

7) Weekly Structure (Sessions 1–6; Intent Per Day; High-Risk Exposures Noted)

\begin{itemize}
\item {[\textbf{REQUIRED}: Use the Weekly Structure table in 1.E.2.3; keep intent-level only; no prescriptions.]}
\item {[\textbf{REQUIRED}: Each session must name Cardio modality intent; do not embed minimum duration in the table header.]}
\end{itemize}

8) Progression Rules (Volume/Intensity Caps; Regression Triggers)

\begin{itemize}
\item Caps: [\textbf{REQUIRED}: Define caps qualitatively (e.g., “no sudden increases,” “stepwise progression,” “impact-managed”). No numeric ramps here.]
\item Regression triggers: [\textbf{REQUIRED}: Trigger/signal language; mechanics-altering pain; repeated invalidity; environment hazards.]
\item One-lever rule: [\textbf{REQUIRED}: State that only one progression lever is adjusted at a time unless safety requires regression.]
\end{itemize}

9) Deload And Test Placement (Mid-Phase And End-Phase)

\begin{itemize}
\item {[\textbf{REQUIRED}: State that mid-phase and end-phase checkpoints exist; specify “mid” and “end” placement conceptually without calendar specifics.]}
\item {[\textbf{REQUIRED}: Reslot posture when safety/validity compromised (tests under compromised conditions are inadmissible).]}
\end{itemize}

10) Test Battery (IDs/Labels Only; Protocols Live Elsewhere)

\begin{itemize}
\item Mid-Phase Tests: [\textbf{REQUIRED}: IDs/names only; no protocol text.]
\item End-Phase Tests: [\textbf{REQUIRED}: IDs/names only; no protocol text.]
\item {[\textbf{REQUIRED}: Use the Test Battery table in 1.E.2.3.]}
\end{itemize}

11) Exit Standards (Objective; Pass/Fail)

\begin{itemize}
\item {[\textbf{REQUIRED}: Objective, auditable exit standards aligned to the phase’s purpose; pass/fail.]}
\item {[\textbf{REQUIRED}: Explicit “cannot pass if” disqualifiers (e.g., unresolved gait-altering pain; evidence inadmissible due to missing logs).]}
\end{itemize}

12) Risk Controls And Stop Rules (Phase-Specific Overlays; Must Reference Stage 0.5)

\begin{itemize}
\item {[\textbf{REQUIRED}: List risk controls; reference \textbf{RISK-1B}-\#\#\# where applicable.]}
\item {[\textbf{REQUIRED}: Stop/abort posture references Stage 0.5 categories; no medical advice.]}
\item {[\textbf{REQUIRED}: Water governance statement if any aquatic exposure is present: certified supervision required for any underwater/breath-hold exposure; unsupervised exposure prohibited.]}
\end{itemize}

13) Required Capabilities And Dependencies (Equipment/Environment/Logging Prerequisites)

\begin{itemize}
\item {[\textbf{REQUIRED}: Reference \textbf{DEP-1C}-\#\#\# / \textbf{ASM-1C}-\#\#\# in the Dependencies table; include fail actions.]}
\item {[\textbf{REQUIRED}: Declare minimum infrastructure dependencies when applicable (footwear interface, timing method, hydration access, safe route/indoor substitute).]}
\end{itemize}

14) Validity Requirements (What Makes Evidence Acceptable)

\begin{itemize}
\item {[\textbf{REQUIRED}: Define admissibility requirements at high level: complete logs, environment fields when applicable, no boundary violations.]}
\item {[\textbf{REQUIRED}: State that missing required documentation makes evidence inadmissible for gates.]}
\item {[\textbf{REQUIRED}: State that Modified/Invalid handling is governed elsewhere; this section may not redefine tags.]}
\end{itemize}

15) Change Control Notes (Freeze Windows; Patch Discipline)

\begin{itemize}
\item {[\textbf{REQUIRED}: State that execution-affecting changes require logged change control and patching; freeze window discipline applies.]}
\item {[\textbf{REQUIRED}: Reference that schema changes require \textbf{PATCH}-1E-\#\#\#; artifact-level changes use the program’s patch conventions elsewhere.]}
\item {[\textbf{REQUIRED}: State that any upstream constraint change triggers downstream revalidation posture.]}
\end{itemize}

16) Compliance Checklist (Pass/Fail)

\begin{itemize}
\item {[\textbf{REQUIRED}: Pass/fail items that enforce template completeness, constraints, dependency references, and evidence posture.]}
\item {[\textbf{REQUIRED}: Must include a fail item for any missing required section, missing required table, or altered column header.]}
\end{itemize}

\newpage
\subsubsection*{1.E.2.3 Required Tables}
\phantomsection

1) Weekly Structure Table

\begingroup
\scriptsize
\setlength{\tabcolsep}{2pt}
\renewcommand{\arraystretch}{1.12}
\begin{longtblr}{
  width=\linewidth,
  rowhead=1,
  colspec = {
    Q[l,wd=0.70in]
    X[l,1.15]
    X[l,1.15]
    X[l,1.30]
    X[l,1.60]
  },
  hlines,
  vlines,
  cells = {hsep=6pt, vsep=6pt, halign=l, valign=t},
  row{odd} = {bg=white},
  row{even} = {bg=LightGray},
  row{1} = {bg=StageBlue!15, fg=black, font=\bfseries, halign=l, valign=t},
}
\Hdr{Session \#} & \Hdr{Primary Focus} & \Hdr{Main Work Domains} & \Hdr{Cardio Modality/Intent} & \Hdr{Notes/Risk Controls} \\
1 & {[\textbf{REQUIRED}]} & {[\textbf{REQUIRED}]} & {[\textbf{REQUIRED}]} & {[\textbf{REQUIRED}]} \\
2 & {[\textbf{REQUIRED}]} & {[\textbf{REQUIRED}]} & {[\textbf{REQUIRED}]} & {[\textbf{REQUIRED}]} \\
3 & {[\textbf{REQUIRED}]} & {[\textbf{REQUIRED}]} & {[\textbf{REQUIRED}]} & {[\textbf{REQUIRED}]} \\
4 & {[\textbf{REQUIRED}]} & {[\textbf{REQUIRED}]} & {[\textbf{REQUIRED}]} & {[\textbf{REQUIRED}]} \\
5 & {[\textbf{REQUIRED}]} & {[\textbf{REQUIRED}]} & {[\textbf{REQUIRED}]} & {[\textbf{REQUIRED}]} \\
6 & {[\textbf{REQUIRED}]} & {[\textbf{REQUIRED}]} & {[\textbf{REQUIRED}]} & {[\textbf{REQUIRED}]} \\
\end{longtblr}
\endgroup

2) Test Battery Table

\begingroup
\scriptsize
\setlength{\tabcolsep}{2pt}
\renewcommand{\arraystretch}{1.12}
\begin{longtblr}{
  width=\linewidth,
  rowhead=1,
  colspec = {
    Q[l,wd=1.00in]
    X[l,1.15]
    Q[l,wd=1.00in]
    X[l,1.40]
    X[l,1.35]
  },
  hlines,
  vlines,
  cells = {hsep=6pt, vsep=6pt, halign=l, valign=t},
  row{odd} = {bg=white},
  row{even} = {bg=LightGray},
  row{1} = {bg=StageBlue!15, fg=black, font=\bfseries, halign=l, valign=t},
}
\Hdr{Test Timing} & \Hdr{Test \textbf{ID}/Name} & \Hdr{Domain} & \Hdr{Validity Conditions} & \Hdr{Pass Threshold} \\
Mid-Phase & {[\textbf{REQUIRED}: \textbf{ID}/Name only]} & {[\textbf{REQUIRED}]} &
\SetCell{halign=l,valign=t}{[\textbf{REQUIRED}: high-level admissibility conditions]} &
\SetCell{halign=l,valign=t}{[\textbf{REQUIRED}: reference threshold source; no protocol text]} \\
End-Phase & {[\textbf{REQUIRED}: \textbf{ID}/Name only]} & {[\textbf{REQUIRED}]} &
\SetCell{halign=l,valign=t}{[\textbf{REQUIRED}: high-level admissibility conditions]} &
\SetCell{halign=l,valign=t}{[\textbf{REQUIRED}: reference threshold source; no protocol text]} \\
\end{longtblr}
\endgroup

3) Dependencies Table

\begingroup
\scriptsize
\setlength{\tabcolsep}{2pt}
\renewcommand{\arraystretch}{1.12}
\begin{longtblr}{
  width=\linewidth,
  rowhead=1,
  colspec = {
    Q[l,wd=1.15in]
    X[l,1.35]
    Q[l,wd=1.35in]
    Q[l,wd=1.55in]
    X[l,1.70]
  },
  hlines,
  vlines,
  cells = {hsep=6pt, vsep=6pt, halign=l, valign=t},
  row{odd} = {bg=white},
  row{even} = {bg=LightGray},
  row{1} = {bg=StageBlue!15, fg=black, font=\bfseries, halign=l, valign=t},
}
\Hdr{Dependency \textbf{ID}} & \Hdr{Required Capability} & \Hdr{Due-By} & \Hdr{Owner} & \Hdr{Fail Action} \\
{[\textbf{DEP-1C}-\#\#\# / \textbf{ASM-1C}-\#\#\#]} & {[\textbf{REQUIRED}]} & {[Phase/Week or event-based]} &
\SetCell{halign=l,valign=t}{[Coach Lead / Trainee / \\ Medical / Equipment Lead]} &
\SetCell{halign=l,valign=t}{[Substitute / Delay escalation / Hold / Regress / Medical Hold + brief action]} \\
\end{longtblr}
\endgroup

\newpage
\subsubsection*{1.E.2.4 Field Definitions}
\phantomsection

A) Header Fields (Phase Header Block)

Phase \textbf{ID} (Required)

\begin{itemize}
\item Allowed values: two-digit 00–13 only. Leading zero required for Phase 00.
\item Completion rule: must match the phase numbering used throughout the binder; must not be renumbered.
\item Status impact if missing/invalid: Draft — Non-Deployable.
\end{itemize}

Phase Name (Required)

\begin{itemize}
\item Allowed values: free text.
\item Completion rule: must remain stable; any meaning change requires change control.
\item Status impact if missing: Draft — Non-Deployable.
\end{itemize}

Duration (weeks) (Required)

\begin{itemize}
\item Allowed values: positive integer.
\item Completion rule: must be provided as a single integer; must not embed calendar dates.
\item If unknown: must remain Draft — Non-Deployable until defined.
\item Status impact if missing/invalid: Draft — Non-Deployable.
\end{itemize}

Version (Required)

\begin{itemize}
\item Allowed values: vX.Y where X and Y are integers; baseline v1.0 unless patched.
\item Completion rule: must reflect the artifact’s controlled state; no silent edits.
\item Status impact if missing/invalid: Draft — Non-Deployable.
\end{itemize}

Governing Status (Required; locked strings)

\begin{itemize}
\item Allowed values only: Draft — Non-Deployable \,|\, Deployable.
\item Definition (must be used verbatim in meaning): Deployable means the artifact’s required fields are complete and its applicable gate checklist has been passed and recorded. Do not append additional label text.
\item Completion rule: default to Draft — Non-Deployable; set to Deployable only after the Phase Row compliance checklist passes and any applicable gate record exists.
\item Status impact if missing/invalid: Draft — Non-Deployable.
\end{itemize}

Owner (primary executor) (Optional)

\begin{itemize}
\item Allowed values: Coach Lead \,|\, Trainee \,|\, Medical \,|\, Equipment Lead.
\item Completion rule: if populated, must match exactly.
\item Status impact: none.
\end{itemize}

Patch References (Optional)

\begin{itemize}
\item Allowed values: \textbf{PATCH}-1E-\#\#\# for schema changes; other patch IDs permitted only if they exist elsewhere and are clearly scoped.
\item Status impact if misleading: Draft — Non-Deployable until corrected.
\end{itemize}

B) Section Completion Standards (Minimums)

All 16 headings are required and must be present exactly as written. Empty sections fail.

1) Phase Mission

\begin{itemize}
\item Must define what the phase accomplishes and what it protects (safety, validity, feasibility).
\item Must include evidence posture statement (decision-grade, admissible).
\item Prohibited: motivational language; programming prescriptions.
\end{itemize}

2) Phase Purpose (Why This Phase Exists)

\begin{itemize}
\item Must state why the phase exists in sequence and what is explicitly not attempted.
\item Prohibited: protocols and calendars.
\end{itemize}

3) Entry Criteria (Gates To Start)

\begin{itemize}
\item Must be objective and auditable; must reference \textbf{DEP}/\textbf{ASM} IDs when dependencies apply.
\item Must include fail posture (Substitute/Delay/Hold/Regress/Medical Hold).
\item Prohibited: subjective readiness language.
\end{itemize}

4) Operating Constraints

\begin{itemize}
\item Must include in body text (required minimum set):
  \begin{itemize}
  \item 6 sessions/week,
  \item required session elements,
  \item Cardio minimum standard (30 minutes continuous per training session) stated in body text,
  \item 10,000 steps/day baseline stated in body text and stated as not a session,
  \item safety override posture,
  \item admissibility posture,
  \item no two-a-days, no make-up doubling,
  \item same-day logging posture.
  \end{itemize}
\item Prohibited: alternate cadence rules.
\end{itemize}

5) Primary Adaptations And Physiological Focus

\begin{itemize}
\item Must remain domain-level; must identify adaptation intent without prescribing methods.
\item Prohibited: set/rep/pace/ramp details.
\end{itemize}

6) Primary Modalities (Strength/Conditioning/Ruck/Swim/Etc.)

\begin{itemize}
\item Must list allowable modality domains and prohibited domains/conditions for the phase.
\item Must respect water governance: underwater/breath-hold unsupervised is always prohibited.
\item Prohibited: underwater/breath-hold instruction.
\end{itemize}

7) Weekly Structure (Sessions 1–6; Intent Per Day; High-Risk Exposures Noted)

\begin{itemize}
\item Weekly Structure table must contain exactly six rows and Session \# must be 1–6.
\item Cardio Modality/Intent must be present for each session row.
\item Prohibited: treating any day as optional.
\end{itemize}

8) Progression Rules (Volume/Intensity Caps; Regression Triggers)

\begin{itemize}
\item Must define caps and regression triggers as auditable rules tied to risk signals (mechanics-altering pain, repeated invalidity, etc.).
\item Prohibited: numeric progression ramps.
\end{itemize}

9) Deload And Test Placement (Mid-Phase And End-Phase)

\begin{itemize}
\item Must include mid-phase and end-phase posture without calendar schedules.
\item Must include validity posture: compromised tests are inadmissible.
\item Prohibited: dates, week-by-week calendars.
\end{itemize}

10) Test Battery (IDs/Labels Only; Protocols Live Elsewhere)

\begin{itemize}
\item Must list IDs/names only; no protocol text.
\item Domain must be specified for each test.
\item Prohibited: step-by-step test instructions.
\end{itemize}

11) Exit Standards (Objective; Pass/Fail)

\begin{itemize}
\item Must be objective and auditable; must include disqualifiers relevant to safety/validity.
\item Prohibited: “best effort” standards.
\end{itemize}

12) Risk Controls And Stop Rules (Phase-Specific Overlays; Must Reference Stage 0.5)

\begin{itemize}
\item Must reference \textbf{RISK-1B}-\#\#\# where applicable.
\item Must explicitly align to Stage 0.5 stop-rule posture; must not weaken boundaries.
\item Prohibited: medical advice.
\end{itemize}

13) Required Capabilities And Dependencies (Equipment/Environment/Logging Prerequisites)

\begin{itemize}
\item Dependencies must use \textbf{DEP-1C}-\#\#\# and \textbf{ASM-1C}-\#\#\# only.
\item Must include fail action outcomes (Substitute/Delay/Hold/Regress/Medical Hold).
\item Prohibited: new dependency \textbf{ID} formats.
\end{itemize}

14) Validity Requirements (What Makes Evidence Acceptable)

\begin{itemize}
\item Must state admissibility posture at governance level: complete logs, environment fields when applicable, no boundary violations.
\item Must state that missing required documentation degrades admissibility.
\item Must not redefine the global meaning of Valid/Modified/Invalid.
\end{itemize}

15) Change Control Notes (Freeze Windows; Patch Discipline)

\begin{itemize}
\item Must state freeze-window posture and patch discipline.
\item Must state that schema changes require \textbf{PATCH}-1E-\#\#\# and that downstream compatibility review is required.
\item Prohibited: silent edits.
\end{itemize}

16) Compliance Checklist (Pass/Fail)

\begin{itemize}
\item Must be explicit pass/fail.
\item Must include fail items for missing section, missing table, altered column headers, and missing operating constraints content.
\item Prohibited: vague checklists.
\end{itemize}

C) Required Table Column Definitions

Weekly Structure Table

\begin{itemize}
\item Session \#: must be 1–6 only.
\item Primary Focus: intent-level descriptor; no prescriptions.
\item Main Work Domains: domain-level labels only.
\item Cardio Modality/Intent: must include Cardio modality category and intent; must not embed minimum duration in the header.
\item Notes/Risk Controls: include risk controls, substitutions, and references to \textbf{RISK-1B}-\#\#\# when applicable.
\end{itemize}

Test Battery Table

\begin{itemize}
\item Test Timing: must be Mid-Phase or End-Phase (unless an additional timing label is explicitly adopted and used consistently).
\item Test \textbf{ID}/Name: IDs/names only; must reference the Stage 2 test inventory when available; no protocol text.
\item Domain: must be a stable domain label; if a new domain label is used, it must be stabilized in Stage 2 and then used consistently.
\item Validity Conditions: high-level admissibility conditions (measured course, unit clarity, environment fields recorded, supervision status if aquatic).
\item Pass Threshold: objective threshold statement or reference to threshold source label; no protocol content.
\end{itemize}

Dependencies Table

\begin{itemize}
\item Dependency \textbf{ID}: \textbf{DEP-1C}-\#\#\# or \textbf{ASM-1C}-\#\#\# only.
\item Required Capability: capability category (not products).
\item Due-By: phase/week or event-based (no calendar dates required).
\item Owner: must be one of Coach Lead | Trainee | Medical | Equipment Lead.
\item Fail Action: must include Substitute, Delay escalation, Hold, Regress, or Medical Hold posture and brief action.
\end{itemize}

D) Conservative Defaults and Placeholder Rules (Binding)

\begin{itemize}
\item Placeholders are permitted only while the Phase Row remains Draft — Non-Deployable.
\item Deployable requires removal of placeholders and completion to minimum standards.
\item Allowed placeholders (Draft only): [\textbf{REQUIRED}], [\textbf{OPTIONAL}], [\textbf{TBD}].
\item Any unresolved required field at the time of attempted deployable labeling forces Draft — Non-Deployable.
\end{itemize}

\newpage
\subsubsection*{1.E.2.5 Phase Row Template Compliance Checklist}
\phantomsection

\textbf{PASS} only if every item below is true:

\begin{itemize}
\item Phase \textbf{ID} format is correct (00–13) and Phase Name is present.
\item Duration (weeks) is present as a positive integer and Version is present as vX.Y.
\item Governing Status is set to an allowed value (Draft — Non-Deployable or Deployable) and is consistent with checklist outcome.
\item All 16 required headings are present exactly as written (no renaming) and meet minimum completion standards.
\item All 3 required tables are present and use exact column headers.
\item Entry Criteria and Exit Standards are objective and pass/fail.
\item Test Battery lists IDs/names only; no protocol text appears anywhere in the Phase Row.
\item Dependencies reference valid \textbf{DEP-1C}-\#\#\# or \textbf{ASM-1C}-\#\#\# IDs, include Due-By, Owner, and explicit Fail Action posture.
\item Operating Constraints section body text includes:
  \begin{itemize}
  \item 6 sessions per week,
  \item required session elements,
  \item Cardio minimum standard stated in body text,
  \item 10,000 steps/day baseline stated in body text and stated as not a session,
  \item safety override posture,
  \item admissibility posture,
  \item no two-a-days / no make-up doubling,
  \item same-day logging posture.
  \end{itemize}
\item Risk Controls section references Stage 0.5 stop-rule posture and references \textbf{RISK-1B}-\#\#\# where applicable.
\item No prohibited content appears (no workouts, no protocols, no calendars, no brand/retailer links, no medical advice).
\end{itemize}

\textbf{FAIL} if any item above is not satisfied:

\begin{itemize}
\item If any required item fails → Governing Status \textbf{MUST} be Draft — Non-Deployable.
\end{itemize}

\subsection*{1.E.3 Conflict-Resolution Rule}
\phantomsection
\addcontentsline{toc}{subsection}{1.E.3 Conflict-Resolution Rule}

\begin{itemize}
\item If any downstream artifact deviates from these templates, omits required fields, changes required headings, changes required table column headers, uses non-canonical IDs or naming, or uses a non-allowed Governing Status value, it \textbf{MUST} be labeled Draft — Non-Deployable until corrected and re-versioned.

\item If a template schema change is required, it \textbf{MUST} be executed via \textbf{PATCH}-1E-\#\#\# with an explicit change-log entry. After a schema patch:
  \begin{itemize}
  \item all instantiated artifacts that use Stage 1.E templates \textbf{MUST} be reviewed for compatibility,
  \item compatible artifacts \textbf{MUST} be updated to the patched schema and retain truthful versioning,
  \item incompatible or unupdated artifacts \textbf{MUST} be labeled Draft — Non-Deployable until brought into compliance.
  \end{itemize}
\end{itemize}

\end{document}